\begin{figure}[htbp]
\centering
\begin{tikzpicture}[font=\small,>=Stealth]
\node at (1,1.5) {interval};
\draw[thick] (0,0)--(2,0);
\fill[purple!80!black] (0,0) circle(.065) node[below,text=black] {$0$};
\fill[purple!80!black] (2,0) circle(.065) node[below,text=black] {$1$};
\draw[->,thick] (2.5,0)--(3.2,0);
\begin{scope}[shift={(4.4,0)}]
\node at (0,1.5) {bend; join endpoints};
\draw[thick] (30:.85) arc(30:330:.85);
\foreach \a in {30,330}{\fill[purple!80!black] (\a:.85) circle(.065);}
\draw[->,purple!80!black] (.85,.4)--(.85,.08);
\draw[->,purple!80!black] (.85,-.4)--(.85,-.08);
\end{scope}
\draw[->,thick] (5.65,0)--(6.35,0);
\begin{scope}[shift={(7.6,0)}]
\node at (0,1.5) {circle};
\draw[thick] (0,0) circle(.85);
\fill[purple!80!black] (.85,0) circle(.065) node[right,text=black] {$[0]=[1]$};
\end{scope}
\end{tikzpicture}
\caption{Joining the interval's two purple endpoints gives one circle
point. A neighborhood of that point joins short intervals near both
endpoints. This explains the seam without mistaking the whole closed
interval for a single circle chart.}
\label{fig:circle-gluing}
\end{figure}
